% Generated draft: requires mathematical and source review.
\documentclass[11pt]{article}
\usepackage[margin=25mm]{geometry}
\usepackage{fontspec}
\setmainfont{Latin Modern Roman}
\usepackage{amsmath,amssymb,mathtools,unicode-math}
\setmathfont{Latin Modern Math}
\usepackage{enumitem,needspace,xurl,hyperref}
\setlength{\parindent}{0pt}
\setlength{\parskip}{5pt}
\setlength{\emergencystretch}{3em}
\newcommand{\N}{\mathbb{N}}
\newcommand{\Z}{\mathbb{Z}}
\newcommand{\Q}{\mathbb{Q}}
\newcommand{\R}{\mathbb{R}}
\newcommand{\C}{\mathbb{C}}
\newcommand{\F}{\mathbb{F}}
\newcommand{\A}{\mathbb{A}}
\newcommand{\PP}{\mathbb{P}}
\newcommand{\E}{\mathbb{E}}
\newcommand{\eps}{\varepsilon}
\newcommand{\dd}{\,\mathrm{d}}
\begin{document}
\section*{UnsolvedMath ID: 30006276}
\textbf{Source problem code:} catalogue entry
\subsection{Definitions and mathematical statement}
Let \(\mathcal A_g\) denote the moduli space of principally polarized abelian varieties of dimension \(g\), as described in the supplied record. Write \(\mathrm{CH}^{*}(\mathcal A_g)\) for its graded Chow ring, with intersection product. The base, coefficient ring, stack versus coarse-space convention, and intended range of \(g\) are not specified in the supplied excerpt.

Let \(T_g\subseteq\mathrm{CH}^{*}(\mathcal A_g)\) denote the subspace called tautological in the source, and let
\[
\pi_g:\mathrm{CH}^{*}(\mathcal A_g)\longrightarrow T_g
\]
denote the tautological projection referred to as equation \((1)\) there. The supplied record does not include that equation, a definition of \(T_g\), or the construction characterizing \(\pi_g\). Thus these symbols refer to the intended source construction; they do not specify a uniquely determined projection from the supplied evidence alone. A projection onto a subspace is understood to restrict to the identity on that subspace.

\textbf{Question, conditional on the missing source definitions.} With the inherited intersection product on \(T_g\), is \(\pi_g\) a ring homomorphism? Its multiplicativity condition is
\[
\pi_g(\alpha\cdot\beta)=\pi_g(\alpha)\cdot\pi_g(\beta)
\qquad\text{for all }\alpha,\beta\in\mathrm{CH}^{*}(\mathcal A_g).
\]
Calling this a ring homomorphism presupposes that \(T_g\) is closed under the inherited product and that the projection has the appropriate additive structure. If ring homomorphisms are required to preserve the identity, the condition also includes \(\pi_g(1)=1\); this follows from the projection property if \(1\in T_g\). These structural conventions require confirmation from the omitted construction. The record also does not specify whether the question concerns a fixed \(g\) or every \(g\) in an intended range.

The surrounding excerpt describes joint work with R. Pandharipande and H.-H. Tseng on maps \(C\to E\), where \(C\) has genus \(g\), \(E\) is an elliptic curve, and both curves vary in moduli. It mentions three approaches connected to three open problems: intersection theory on \(\mathcal A_g\), Gromov--Witten theory of the universal elliptic curve \(\mathcal E\to\mathcal M_{1,1}\), and genus-one Gromov--Witten theory of the Hilbert scheme of points in \(\C^2\). This context supplies no further definition of the projection in the present record.
\subsection{Short English statement}
The question asks whether projecting Chow classes on \(\mathcal A_g\) onto the source's tautological subspace preserves multiplication. The defining projection and several domain conventions are omitted, so this is a conditional statement requiring the original construction. The supplied dated curation labels the problem partially solved; that assessment is not independently verified here.
\subsection{Sources}
\begin{itemize}
\item Supplied local UnsolvedMath registry record, permanent ID \(30006276\), title \emph{Multiplicativity of Tautological Chow Projections}, exact source problem code \texttt{catalogue entry}. The record gives proposed year 2025 and identifies no proposer.
\item \emph{Recent Trends in Algebraic Geometry} (2025), Oberwolfach Reports, pp.\ 1451--1516, DOI \texttt{10.4171/owr/2025/28}. Supplied report reference: \url{https://doi.org/10.4171/owr/2025/28}. The excerpt calls this Question 1, refers to a projection in equation \((1)\), and says that the question was introduced in reference [3]. Neither equation \((1)\) nor the bibliographic identity of [3] is supplied.
\item Supplied literature source: Iribar L\'opez, \emph{Noether--Lefschetz cycles}, \url{https://doi.org/10.1017/fmp.2025.10018}.
\item Supplied literature source labelled ``original OWR report'': \url{https://doi.org/10.4171/owr/2025/28}.
\item Literature assessment dated exactly \texttt{2026-08-22}: the supplied curation labels the problem partially solved and states that the general multiplicativity question remains unresolved, while Iribar L\'opez proves the homomorphism property in several cases and for a natural subring generated by Noether--Lefschetz data. This is supplied dated curation, not an independently verified literature assessment; the record itself requests independent verification of time-sensitive status and cited claims. The cases and subring are not specified in the supplied evidence.
\item The registry labels the statement \texttt{corrected\_verified} and reports that the official Oberwolfach report was checked, the question isolated, and its notation repaired. These are supplied source-verification claims, not verification performed for this draft. The surrounding excerpt attributes the enumerative work to joint work with R. Pandharipande and H.-H. Tseng without identifying the speaker in the supplied text.
\end{itemize}
\end{document}
